\section{Perturbation-only Mandelbrot rendering (without linear approximation)}
\label{sec:perturb-only}

This kernel can render the Mandelbrot set using \emph{perturbation alone}, i.e.,
without taking any LA~v2 linear-approximation steps. The same CUDA entry point
\code{mandel\_1xHDR\_float\_perturb\_lav2<IterType,T,SubType,Mode,PExtras>} is
used; the behavior is selected at compile time via \code{LAv2Mode}. In
particular, when \code{Mode} includes only the perturbation path (e.g.\
\code{LAv2Mode::PO}), the kernel skips the LA stage traversal and runs only the
perturbation loop against a stored reference orbit (\cref{sec:ref-orbit-calc}).

\subsection{Rendering objective}
\label{sec:perturb-only-goal}

The goal remains standard escape-time rendering for
\[
z_{n+1} = z_n^2 + c,\qquad z_0=0,
\]
with bailout \(|z|^2 \ge 4\). For each pixel, the kernel computes the parameter
\(c\), iterates until escape or \code{n\_iterations}, and stores the resulting
iteration count in \code{OutputIterMatrix[idx]}.

In perturbation rendering, the expensive high-precision orbit evaluation for
each pixel is avoided by reusing a \emph{reference orbit} computed at a
reference parameter \(c_\star\), then evolving only the \emph{delta orbit} for
nearby pixels.

\subsection{Reference orbit and delta formulation}
\label{sec:perturb-only-deltas}

For a pixel parameter \(c\) near the reference parameter \(c_\star\), define
\(\Delta c = c - c_\star\) and \(\Delta z_n = z_n - z_n^\star\).
Substituting into the Mandelbrot recurrence yields the exact delta recurrence
derived in \cref{subsec:perturbation-mixed-precision}:
\cref{eq:perturb-exact}.
Perturbation uses this recurrence directly: it is not a linearization. The work
per iteration is reduced because \(z_n^\star\) is fetched from storage rather
than recomputed in high precision.

\subsection{Pixel parameter delta \texorpdfstring{$\Delta c$}{Delta c}}
\label{sec:perturb-only-deltac}

For each pixel \((X,Y)\), the kernel constructs a delta parameter relative to a
chosen reference center (sign conventions incorporate the image mapping):
\begin{align}
\Delta c_x &= dx \cdot X - \code{centerX}, \\
\Delta c_y &= -dy \cdot Y - \code{centerY},
\end{align}
and stores this as \code{DeltaSub0} (with scalar components \code{DeltaSub0X},
\code{DeltaSub0Y}). The perturbation state starts at
\[
\Delta z_0 = 0,
\]
so \code{DeltaSubN} is initialized to zero.

\subsection{Perturbation recurrence used in the kernel}
\label{sec:perturb-only-update}

Writing the reference sample as \(z^\star = x^\star + i y^\star\) and the delta
as \(\Delta z = \Delta x + i \Delta y\), the delta recurrence can be expressed
in a factored form convenient for implementation:
\begin{equation}
\Delta z \leftarrow \Delta z \cdot (2 z^\star + \Delta z) + \Delta c.
\label{eq:perturb-factor}
\end{equation}
This identity is equivalent to
\(\Delta z_{n+1} = 2 z_n^\star \Delta z_n + (\Delta z_n)^2 + \Delta c\), because
\(\Delta z\cdot(2z^\star+\Delta z) = 2z^\star\Delta z + (\Delta z)^2\).

The kernel implements \cref{eq:perturb-factor} in real arithmetic by forming
the sums
\[
(2x^\star + \Delta x),\qquad (2y^\star + \Delta y),
\]
then updating:
\begin{align}
\Delta x' &= \Delta x\,(2x^\star + \Delta x) - \Delta y\,(2y^\star + \Delta y) + \Delta c_x, \\
\Delta y' &= \Delta x\,(2y^\star + \Delta y) + \Delta y\,(2x^\star + \Delta x) + \Delta c_y.
\end{align}
Note that both right-hand sides use the \emph{pre-update} values of
\(\Delta x\) and \(\Delta y\).  In the code, these originals are saved as
\code{DeltaSubNXOrig} and \code{DeltaSubNYOrig} before either component is
overwritten.  The intermediate quantities correspond to:
\begin{lstlisting}
tempSum1 = 2*zy + DeltaSubNYOrig;  // (2 y^\star + \Delta y)
tempSum2 = 2*zx + DeltaSubNXOrig;  // (2 x^\star + \Delta x)
\end{lstlisting}
followed by the real/imag updates. For extended or HDR types (\cref{sec:hdr-float}), the kernel routes
the same math through type-specialized implementations
(\code{T::custom\_perturb2} / \code{T::custom\_perturb3}) to keep the inner loop
tight and to enforce the type's reduction/normalization rules.

\subsection{Reconstructing the absolute orbit for escape testing}
\label{sec:perturb-only-reconstruct}

Perturbation evolves \(\Delta z_n\), but escape-time rendering requires a
bailout test on the absolute orbit \(z_n\). Each iteration reconstructs:
\[
z_n \approx z_n^\star + \Delta z_n,
\]
using the stored reference sample \(z_n^\star\) and the current delta. In the
kernel, this appears as:
\[
x = x^\star + \Delta x,\qquad y = y^\star + \Delta y.
\]
The bailout test is then performed on
\[
|z|^2 = x^2 + y^2,
\]
with a squared bailout threshold of~256 in this CUDA entry point.
The loop continues only while $|z|^2<256$, so equality stops evaluation.
The mathematical Mandelbrot escape radius~2 gives squared magnitude~4;
the implemented rendering bailout is later, at radius~16. The newly
reconstructed value is tested before the stored iteration count is
increased for that update. For HDR and related types, the norm and
comparison use reduced values and specialized reduced comparators.

\subsection{Reference index management and rebasing}
\label{sec:perturb-only-rebase}

The perturbation loop advances a reference-orbit index \code{RefIteration} by
one sample per update, fetching \(z^\star\) samples from \code{Perturb}. Rebasing
resets that index while preserving the accumulated \code{iter} count. The
kernel includes a \emph{rebasing} mechanism that resets the delta
representation when it becomes ill-conditioned. Conceptually, if the delta
becomes comparable to or larger than the reconstructed orbit, the decomposition
\(z = z^\star + \Delta z\) stops being numerically advantageous. In that case,
the kernel folds the delta into the base by setting:
\[
\Delta z \leftarrow z,\qquad \text{and restart the reference index.}
\]
Operationally, the code compares the reconstructed orbit magnitude proxy
against the delta magnitude proxy, and also rebasing if the reference index
reaches the end of the stored orbit samples (to avoid out-of-range sampling).
After rebasing, perturbation continues from the new base representation.

This mechanism keeps perturbation theory usable across a wide range of
pixels and iteration depths while preserving the core rendering objective:
compute escape-time using a stable bailout on the reconstructed orbit.

\subsection{Using \code{LAv2Mode} to select perturbation-only execution}
\label{sec:perturb-only-mode}

The kernel is structured as two compile-time phases:
\begin{itemize}
\item an LA~v2 phase guarded by \code{Mode == Full || Mode == LAO},
\item a perturbation phase guarded by \code{Mode == Full || Mode == PO}.
\end{itemize}
Therefore, perturbation-only rendering is achieved by instantiating the kernel
with a mode that includes only the perturbation path (e.g.\ \code{LAv2Mode::PO}).
In this configuration:
\begin{itemize}
\item \code{DeltaSub0} is computed from the pixel location,
\item \code{DeltaSubN} remains initialized to \(\Delta z_0 = 0\),
\item the kernel runs the perturbation loop, reconstructing \(z\) each step for
      bailout tests,
\item the final escape-time count is written to \code{OutputIterMatrix}.
\end{itemize}
This kernel provides a complete Mandelbrot renderer based purely on reference-orbit
perturbation, without relying on any linear-approximation hierarchy.
After all pixels are computed, the antialiasing and reduction passes
(\cref{subsec:antialiasing,subsec:reduction}) extract the final color
image and iteration statistics.

While perturbation theory drastically reduces per-pixel cost, the iteration
count can still be enormous at extreme depths.  The next chapter introduces
the linear approximation framework, which skips large blocks of iterations
entirely by replacing them with a single linear map.
